\section{AC-MOT Method}\label{sec:method}

\subsection{Overview and operating point}
Figure~\ref{fig:pipeline} shows the controlled pipeline. Each frame passes
through scene analysis, the AC-MOT controller, the detector, and the tracker.
The controller chooses the detector operating point
\begin{equation}
\theta_t = (c_t, \eta_t, r_t),
\label{eq:theta}
\end{equation}
where $c_t$ is the detector confidence threshold, $\eta_t$ is the NMS IoU
threshold, and $r_t$ is the input resolution (the long side of the resized
image, in pixels). The detector resizes the frame to $r_t$, suppresses boxes
that overlap a higher-scoring box by more than $\eta_t$, and removes boxes
with a score below $c_t$. The remaining boxes go to the tracker. The detector
weights and the tracker are frozen; AC-MOT only changes $\theta_t$.

The controller uses only information that is available at run time. Its
decision for frame $t$ is
\begin{equation}
\theta_t = \pi\left(I_t, \mathcal{B}_{1:t-1}\right),
\label{eq:policy}
\end{equation}
where $I_t$ is the current image and $\mathcal{B}_{1:t-1}$ are the boxes that
the tracker reported for earlier frames. The current image is available when
the decision is made, so its statistics may be used. Tracker output of frame
$t$ is stored only after detection and association, and it is used from frame
$t+1$ on. No ground truth and no future frame enter the decision. The scene
layer reads box coordinates and image statistics only; detector scores do not
enter it.

\begin{figure*}[t]
\centering
\begin{tikzpicture}[
  font=\footnotesize,
  node distance=4mm,
  stage/.style={draw, rounded corners=2pt, align=center, text width=20mm,
                minimum height=11mm, inner sep=2pt},
  ctrl/.style={stage, fill=orange!10},
  host/.style={stage, fill=blue!7},
  note/.style={align=center, text width=22.5mm, font=\footnotesize, inner sep=1pt},
  arr/.style={-{Stealth[length=2mm]}, thick},
  fb/.style={-{Stealth[length=2mm]}, thick, dashed, green!40!black}]
\node[stage] (frame) {Frame $I_t$};
\node[ctrl, right=of frame] (scene) {Scene analysis};
\node[ctrl, right=of scene] (ctrl) {AC-MOT\\controller};
\node[host, right=of ctrl] (det) {Detector\\(frozen)};
\node[host, right=of det] (trk) {Tracker\\(frozen)};
\node[stage, right=of trk] (out) {Tracks of\\frame $t$};
\draw[arr] (frame) -- (scene);
\draw[arr] (scene) -- (ctrl);
\draw[arr] (ctrl) -- node[above, font=\footnotesize] {$\theta_t$} (det);
\draw[arr] (det) -- (trk);
\draw[arr] (trk) -- (out);
\node[note, below=2mm of scene] {edges, brightness, blur of $I_t$; crowd and tiny ratio from tracks of $t-1$};
\node[note, below=2mm of ctrl] {SCI, scene label, recovery, hysteresis, dwell, state};
\node[note, below=2mm of det] {resolution $r_t$, NMS IoU $\eta_t$, score threshold $c_t$};
\node[note, below=2mm of trk] {ByteTrack or OATrack, unchanged};
\draw[fb] (trk.north) -- ++(0,5mm) -| node[pos=0.25, above, font=\footnotesize]
  {tracked boxes, used from frame $t+1$} (scene.north);
\end{tikzpicture}

\caption{AC-MOT in a tracking-by-detection pipeline. Scene analysis and the
controller choose the operating point $\theta_t=(c_t,\eta_t,r_t)$ before the
detector runs. The frozen detector and tracker are not modified. Tracked boxes
of frame $t$ feed the scene analysis of later frames (dashed arrow).}
\label{fig:pipeline}
\end{figure*}

\subsection{Scene analysis}
Scene analysis runs on frame 1 and then every ten frames ($t=1$ or
$t \bmod 10 = 1$); analysis step $j$ denotes the $j$-th such frame. Two kinds
of cues are computed.

\emph{Image cues of the current frame.} The frame is reduced to quarter scale
and converted to gray. The edge cue $e_t$ is the fraction of Canny edge pixels
(thresholds 50 and 120), the brightness cue $\beta_t$ is the mean gray level,
and the blur cue $\nu_t$ is the variance of the Laplacian.

\emph{Box cues of earlier frames.} Let $N_j$ be the number of boxes that the
tracker reported for frame $t-1$. The crowd cue and the tiny-object ratio are
\begin{align}
C_j&=\min\!\left(\frac{N_j}{30},1\right),\label{eq:crowd}\\
T_j&=\frac{1}{N_j}\sum_{i=1}^{N_j}\mathbf{1}\!\left[a_i<32^2\right],\label{eq:tiny}
\end{align}
where $a_i$ is the area of box $i$ in image pixels and $T_j=0$ when $N_j=0$.

\subsection{Crowd-selected SCI and scene label}
The SCI of the frozen controller is the mean of the most recent crowd values,
\begin{equation}
S_j=\frac{1}{|H_j|}\sum_{k\in H_j} C_k,
\label{eq:sci}
\end{equation}
where $H_j$ holds the last seven analysis steps (fewer at the start of a
sequence). The weighted SCI uses the crowd cue alone because the cue audit of
Sect.~\ref{sec:modern} selected the subset \{crowd\}. The other cues are still
computed and still act on the state, as guard and auxiliary cues. They define
a scene label $\ell_j$:
\emph{night} if $\beta_t<80$; otherwise \emph{blur} if $\nu_t<180$; otherwise
\emph{tiny} if $T_j>0.50$; otherwise \emph{crowded} if $C_j>0.65$ or
$e_t>0.13$; otherwise \emph{clear}. The tiny ratio also guards the HIGH state
directly, as shown next.

\subsection{Three-state controller}
The controller has three states, LOW, MEDIUM, and HIGH. The names describe
scene difficulty. They are not a ranking of input resolution: in the frozen
action map the HIGH state uses the smallest input (Table~\ref{tab:actions}).
At each analysis step the state is updated in the order of
Fig.~\ref{fig:state} and Algorithm~\ref{alg:controller}.

\emph{Target.} The target is HIGH if $S_j>\tau_h$ or $T_j>0.50$, MEDIUM if
$S_j>\tau_m$ or $\ell_j\in\{\text{crowded},\text{tiny}\}$, and LOW otherwise,
with the frozen thresholds $\tau_m$ and $\tau_h$ of Table~\ref{tab:constants}.

\emph{Recovery probe.} The controller keeps a decaying peak $P_j$ of the box
count. When $P_{j-1}\geq5$ and $N_j<0.40\,P_{j-1}$, a drop age $d_j$ is
increased; otherwise $d_j=0$. The peak is then updated as
$P_j=\max(N_j,\lfloor 0.8\,P_{j-1}\rfloor)$. While $1\leq d_j\leq3$, the
target is set to HIGH. This gives a short, bounded response when the number
of tracked objects collapses.

\emph{Hysteresis.} Only downward moves are held. In the HIGH state, a lower
target is ignored while $S_j>0.50$ or $T_j>0.40$. In the MEDIUM state, a lower
target is ignored while $S_j>0.25$.

\emph{Dwell.} The state cannot change within 30 frames of the last change.
The controller starts in LOW, and the state is held between analysis frames.

\begin{figure}[t]
\centering
\begin{tikzpicture}[
  font=\footnotesize,
  node distance=3mm,
  blk/.style={draw, rounded corners=2pt, align=left, text width=70mm,
               inner sep=3pt, fill=orange!8},
  arr/.style={-{Stealth[length=2mm]}, thick}]
\node[blk] (cues) {\textbf{1 Cues} (frames $1, 11, 21, \dots$): crowd $C_j$ and
  tiny ratio $T_j$ from the tracked boxes of frame $t-1$; edges, brightness,
  and blur of $I_t$.};
\node[blk, below=of cues] (sci) {\textbf{2 SCI and label}: $S_j$ = mean of the
  last seven $C$ values; label night, blur, tiny, crowded, or clear.};
\node[blk, below=of sci] (target) {\textbf{3 Target}: HIGH if $S_j>\tau_h$ or
  $T_j>0.50$; MEDIUM if $S_j>\tau_m$ or label crowded or tiny; else LOW.};
\node[blk, below=of target] (probe) {\textbf{4 Recovery}: if the box count
  collapses below 0.40 of its decaying peak, target HIGH for up to three
  analyses.};
\node[blk, below=of probe] (hold) {\textbf{5 Hysteresis}: a downward move is
  held while the scene stays difficult.};
\node[blk, below=of hold] (dwell) {\textbf{6 Dwell}: no state change within 30
  frames of the last change.};
\node[blk, below=of dwell, fill=blue!7] (act) {\textbf{7 Action}: the state
  selects $(r_t, \eta_t, c_t)$ from the frozen map (Table~\ref{tab:actions}) and
  is held until the next analysis.};
\foreach \a/\b in {cues/sci, sci/target, target/probe, probe/hold, hold/dwell, dwell/act}
  \draw[arr] (\a) -- (\b);
\end{tikzpicture}

\caption{Update order of the frozen three-state controller at an analysis
frame.}
\label{fig:state}
\end{figure}

\begin{algorithm}[t]
\caption{Frozen AC-MOT controller for one sequence}\label{alg:controller}
\begin{algorithmic}[1]
\State $q\gets$ LOW; $H\gets\emptyset$; $P,d\gets0$; $t_c\gets1$
\For{each frame $t=1,2,\dots$}
  \If{$t=1$ \textbf{or} $t \bmod 10=1$}
    \State $e_t,\beta_t,\nu_t\gets$ statistics of $I_t$
    \State $N,C,T\gets$ box cues of frame $t-1$
    \State add $C$ to $H$ (keep 7); $S\gets\mathrm{mean}(H)$
    \State $\ell\gets$ scene label; $g\gets$ target from $S,T,\ell$
    \If{$P\geq5$ \textbf{and} $N<0.4P$} $d\gets d+1$
    \Else{} $d\gets0$
    \EndIf
    \State $P\gets\max(N,\lfloor0.8P\rfloor)$
    \If{$1\leq d\leq3$} $g\gets$ HIGH
    \EndIf
    \State apply the downward holds of Table~\ref{tab:constants}
    \If{$t-t_c\geq30$ \textbf{and} $g\neq q$}
      \State $q\gets g$; $t_c\gets t$
    \EndIf
  \EndIf
  \State $(r_t,\eta_t,c_t)\gets$ action of state $q$
  \State detect with $(r_t,\eta_t)$; drop scores $<c_t$
  \State update the tracker; store its boxes
\EndFor
\end{algorithmic}
\end{algorithm}

\subsection{Frozen constants and action map}
Table~\ref{tab:constants} lists the frozen constants and
Table~\ref{tab:actions} the frozen action map. LOW uses the largest input
(1536 pixels) with loose NMS and a low score threshold. MEDIUM uses 1280
pixels with tight NMS and a score threshold of 0.40. HIGH uses 1088 pixels
with NMS 0.60 and a score threshold of 0.40. The actions were selected on
calibration data together with $\tau_m$ and $\tau_h$
(Sect.~\ref{sec:modern}). The score threshold $c_t$ is applied on the detector
side before the tracker. It is separate from any confidence gate inside the
tracker, such as the fixed tracker-side gate of OATrack.

\subsection{How the constants were set}\label{sec:constants}
The constants were set with parameter sweeps and experiments, not by hand.
The analysis stride and the SCI window come from a temporal sweep on the
VisDrone validation sequences. It tested 25 combinations, with windows of 1,
3, 5, 7, and 9 analyses and strides of 1, 5, 10, 15, and 20 frames. Among the
settings that met the speed and identity constraints (at least 25 FPS and at
most 271 IDS), a window of 7 and a stride of 10 gave the highest MOTA.
Operating-point sweeps on the same data covered input resolutions from 512 to
960 pixels, confidence thresholds from 0.05 to 0.50, and NMS thresholds from
0.30 to 0.80, and defined the supported values of the first optimized
profiles (Sect.~\ref{sec:development}). The cue subset, the thresholds $\tau_m$
and $\tau_h$, and the action map were selected by cross-validated search on
separate calibration sequences (Sect.~\ref{sec:modern}). The remaining
constants of Table~\ref{tab:constants} were kept fixed in all development
stages.

\begin{table}[t]
\caption{Frozen constants of the AC-MOT scene layer.}
\label{tab:constants}
\centering\footnotesize
\begin{tabular}{@{}p{43mm}l@{}}
\toprule
Constant & Value\\
\midrule
Analysis stride (frames) & 10\\
SCI window (analysis steps) & 7\\
Crowd normalizer (boxes) & 30\\
Tiny-object area (pixels) & $32\times32$\\
MEDIUM threshold $\tau_m$ & 0.29747709701135766\\
HIGH threshold $\tau_h$ & 0.5114928923560124\\
Tiny guard for HIGH & 0.50\\
Label thresholds: brightness, blur & 80, 180\\
Label thresholds: crowd, edges & 0.65, 0.13\\
Recovery: min.\ peak, drop ratio, peak decay, steps & 5, 0.40, 0.80, 3\\
Holds: HIGH ($S$, $T$), MEDIUM ($S$) & (0.50, 0.40), 0.25\\
Dwell (frames) & 30\\
Initial state & LOW\\
\bottomrule
\end{tabular}
\end{table}

\begin{table}[t]
\caption{Frozen action map. State names describe scene difficulty, not a
resolution rank.}
\label{tab:actions}
\centering\footnotesize
\begin{tabular}{@{}lrrr@{}}
\toprule
State & Resolution $r$ & NMS IoU $\eta$ & Score threshold $c$\\
\midrule
LOW & 1536 & 0.70 & 0.10\\
MEDIUM & 1280 & 0.45 & 0.40\\
HIGH & 1088 & 0.60 & 0.40\\
\bottomrule
\end{tabular}
\end{table}
